\documentclass[10pt,a4paper]{amsart}
\usepackage[margin=19mm]{geometry}
\usepackage{amsmath,amssymb,mathtools}
\usepackage[scr=rsfs,cal=euler]{mathalfa}
\usepackage{xfrac}
\usepackage[unicode,hidelinks,bookmarks=false]{hyperref}
\newcommand{\ie}{i.e.\ }
\newcommand{\cf}{cf.\ }
\newcommand{\resp}{respectively\ }
\newcommand{\Subsection}{\subsection}
\providecommand{\nobreakdash}{\nobreak\hbox{-}}
\setlength{\emergencystretch}{2em}
\multlinegap0pt
\usepackage{bm}
\usepackage{bbm}
\usepackage{dsfont}
\usepackage{xspace}
\usepackage{cancel}
\usepackage{mathtools}
\usepackage{amsthm}
\newtheorem{lemma}{Lemma}
\newtheorem{theorem}{Theorem}
\newcommand\argmin{\mathop{\rm argmin}}
\newcommand\bz{\bm{0}}
\newcommand\e{\bm{e}}
\newcommand\mN{\mathcal{N}}
\newcommand\proj{\mathrm{proj}}
\newcommand\q{\bm{q}}
\newcommand\s{\bm{s}}
\newcommand\w{\bm{w}}
\newcommand\y{\bm{y}}
\newcommand\z{\bm{z}}
\newcommand{\R}{{\mathbb{R}}} %Real numbers
\renewcommand\r{\bm{r}}
\renewcommand\u{\bm{u}}
\renewcommand\v{\bm{v}}
\renewcommand{\a}{\bm{a}}
\renewcommand{\c}{\bm{c}}
\renewcommand{\d}{\bm{d}}
\title[Data-Dependent Complexity of First-Order Methods for Support]{Data-Dependent Complexity of First-Order Methods for Support Vector Machines (Lemma 7)}
\author{Matthew Hough, Stephen A. Vavasis}
\date{Source: arXiv:2512.03947v2 (CC BY 4.0)}
\begin{document}
\maketitle

\noindent\emph{Source and licence.} Data-Dependent Complexity of First-Order Methods for Support Vector Machines. Version arXiv:2512.03947v2 (CC BY 4.0), distributed under the Creative Commons Attribution 4.0 International licence, as recorded in the arXiv licence metadata for this exact version and shown on the abstract page https://arxiv.org/abs/2512.03947v2. Full rights notice: licenses/arxiv-2512.03947-CC-BY-4.0.txt.\par\smallskip

\noindent\textit{Editorial scope.} Counted unit: the Theorem printed as Lemma 7 of the article (source0.tex lines 1953--1958, source label \texttt{thm:Csvm-single-point-certificate}), delivered together with the article's own complete original proof of it (source0.tex lines 1959--2034, one \texttt{proof} environment of 76 lines). No mathematics is written, repaired or re-derived: statement and proof are the article's own text line for line. Required context retained uncounted: thm:dist-to-opt-dual (lines 394--405); lem:dual-psi-relation (lines 522--535); prob:sm-svm-dual-p (lines 1792--1797); thm:Csvm-dist-to-opt-dual (lines 1840--1850); thm:Csvm-dual-psi-relation (lines 1926--1939). Every definition, display and label of those lines is retained unchanged. Omitted from this package and not redistributed: 1--393, 406--521, 536--1791, 1798--1839, 1851--1925, 1940--1952, 2035--2291. Nothing inside a retained excerpt is omitted or altered: every retained body is delivered exactly as the article's lines are written, so no omission receipt of any kind accompanies this package. Citations: the retained excerpts carry no citation command at all, so no bibliography entry of the article is transcribed and no reference is redistributed. Reference adaptation: none was needed and none was made; every reference inside the delivered unit resolves inside it, so no unresolved reference or citation is produced. Printed slips kept literal: no printed word, symbol, hypothesis or number of the counted statement or of its proof was altered. Layout and class: the article's own class is \texttt{article} with packages including jmlr2e, amsmath, amssymb, mathtools, tikz, bm, booktabs, algorithm, algpseudocode, lastpage; the wrapper is the lane's pinned \texttt{amsart} editorial wrapper at 10pt on A4, which restates the theorem-like environments the retained text needs and adds the article's own macro definitions that the retained text needs (\texttt{\textbackslash R}, \texttt{\textbackslash a}, \texttt{\textbackslash argmin}, \texttt{\textbackslash bz}, \texttt{\textbackslash c}, \texttt{\textbackslash d}, \texttt{\textbackslash e}, \texttt{\textbackslash mN}, \texttt{\textbackslash proj}, \texttt{\textbackslash q}, \texttt{\textbackslash r}, \texttt{\textbackslash s}, \texttt{\textbackslash u}, \texttt{\textbackslash v}, \texttt{\textbackslash w}, \texttt{\textbackslash y}, \texttt{\textbackslash z}), with extra wrapper packages (bm, bbm, dsfont, xspace, cancel, mathtools). The delivered numbering is the unit's own. Assets: no retained span contains a figure environment, a graphics-inclusion command, a PDF-page inclusion, a table or a TikZ picture; no image file is copied into this package, the package declares no asset dependency and no image file is redistributed with it. Licence: version arXiv:2512.03947v2 of this article is distributed under the Creative Commons Attribution 4.0 International licence, as recorded in the arXiv licence metadata for this exact version and shown on the abstract page https://arxiv.org/abs/2512.03947v2. A full rights notice is delivered at licenses/arxiv-2512.03947-CC-BY-4.0.txt. Characters: every retained excerpt is pure ASCII, so the delivered engine, XeLaTeX with its default Latin Modern text font, reports no missing character.
\begin{theorem}
\label{thm:dist-to-opt-dual}
Let $\y\in\mathbb{R}^n$, and define $\r := \proj_{\Omega}(\y + \frac{1}{L}\nabla g(\y))$, $\epsilon := \kappa \|\y - \r\|$.
Then
\[
    \|\r-\r^*\|\leq\epsilon.
\]
In particular, the FISTA iterates satisfy
\[
    \|\r_{k+1} - \r^*\| \leq \epsilon_{k+1}^{\mathrm{FISTA}}.
\]
\end{theorem}

\begin{lemma} \label{lem:dual-psi-relation}
    Let $(\u^*,\v^*)$ be the unique optimal solution of the perturbed H-SVM dual, and let $(\w^*,t^*)$ be the corresponding optimal classifier. Then
    \[
        u_i^* = \proj_{[0,\gamma]}\left(\frac{1-\theta^*_{\c,i}}{\mu}\right),\qquad \forall i\in[j],
    \]
    and
    \[
        v_i^* = \proj_{[0,\gamma]}\left(\frac{1-\theta^*_{\d,i}}{\mu}\right),\qquad \forall i\in[l].
    \]
    Equivalently,
    \[
        u_i^* = -\psi'(\theta^*_{\c,i}),\quad\forall i \in [j],\quad\text{and}\quad v_i^* = -\psi'(\theta^*_{\d,i}),\quad\forall i \in [l].
    \]
\end{lemma}

\begin{equation} \label{prob:sm-svm-dual-p}
  \begin{array}{rrl}
    &\max&g(\r):=-\frac{1}{2}\Vert H^T\r\Vert^2 + \e^T\r - \frac{\mu}{2}\Vert\r\Vert^2\\
    &\mbox{s.t.}&\r\in \Omega,
  \end{array}
\end{equation}

\begin{theorem} \label{thm:Csvm-dist-to-opt-dual}
Let $\y\in\mathbb{R}^n$, and define $\r := \proj_{\Omega}(\y + \frac{1}{L}\nabla g(\y))$, $\epsilon := \kappa \|\y - \r\|$.
Then
\[
    \|\r-\r^*\|\leq\epsilon.
\]
In particular, the FISTA iterates satisfy
\[
    \|\r_{k+1} - \r^*\| \leq \epsilon_{k+1}^{\mathrm{FISTA}}.
\]
\end{theorem}

\begin{theorem} \label{thm:Csvm-dual-psi-relation}
Let $(\u^*,\v^*)$ be the unique optimal solution of \eqref{prob:sm-svm-dual-p}, and let $(\w^*,t^*)$ be any optimal primal classifier. Then
\[
    u_i^* = \proj_{[0,\gamma]}\left(\frac{1-\c_i^T\w^* - t^*}{\mu}\right),\qquad \forall i\in[j],
\]
and
\[
    v_i^* = \proj_{[0,\gamma]}\left(\frac{1+\d_i^T\w^* + t^*}{\mu}\right),\qquad \forall i\in[l].
\]
Equivalently,
\[
    u_i^* = -\psi'(\c_i^T\w^* + t^*),\quad\forall i \in [j],\quad\text{and}\quad v_i^* = -\psi'(-\d_i^T\w^* - t^*),\quad\forall i \in [l].
\]
\end{theorem}

\begin{theorem} \label{thm:Csvm-single-point-certificate}
Suppose $\mu\gamma<1$. Let $\r = (\u,\v)$ be the projected gradient step for some $\y \in \R^n$, and $\epsilon = \kappa \|\y - \r\|$, as in Theorem~\ref{thm:Csvm-dist-to-opt-dual}. If $\gamma-r_i>\epsilon$, then
\[
    \theta_i(\w(\r),t(\r)) > 1-\mu\gamma > 0, \qquad \theta_i(\w(\r^*),t(\r^*)) > 1-\mu\gamma>0.
\]
\end{theorem}

\begin{proof}
Let $\w := H^T\r$ and $t := t(\r) \in T(\w(\r))$ be chosen arbitrarily. Define $\bar{\r} := (\bar{\u}, \bar{\v})$ by
\begin{align*}
    \bar{u}_i &:= \proj_{[0,\gamma]}\left(\frac{1-\c_i^T\w - t}{\mu}\right),\quad\forall i \in [j],\\
    \bar{v}_i &:= \proj_{[0,\gamma]}\left(\frac{1+\d_i^T\w + t}{\mu}\right),\quad\forall i \in [l].
\end{align*}
We showed in the proof of Lemma~\ref{lem:dual-psi-relation} that
\[
    \psi'(\theta) = -\proj_{[0,\gamma]}\left(\frac{1-\theta}{\mu}\right).
\]
Therefore, using that $t \in \argmin_{s} f(\w,s)$, we have
\begin{align*}
    0 &= \frac{\partial f(\w,t)}{\partial t}\\
      &= \sum_{i=1}^j \psi'(\c_i^T\w + t) - \sum_{i=1}^l \psi'(-\d_i^T\w - t)\\
      &= -\e^T\bar{\u} + \e^T\bar{\v}.
\end{align*}
So $\bar{\r} \in \Omega$. Define the vector $\a$ such that
\[
    a_i = \begin{cases}
        1, & i \in \{1,\hdots,j\},\\
        -1, & i \in \{j+1,\hdots,n\}.
    \end{cases}
\]
We may write
\[
    \bar{\r} = \proj_{[0,\gamma]^n}\left(\frac{\e - H\w - \a t}{\mu}\right),
\]
which has optimality conditions
\[
    \bar{\q} := \e - H\w - \a t - \mu\bar{\r} \in \mN_{[0,\gamma]^n}(\bar{\r}).
\]
Define
\[
    \bar{\z} := \bar{\q} + \a t = \e - H\w - \mu\bar{\r}.
\]
For every $\s \in \Omega$, $\a^T(\s - \bar{\r}) = 0$, therefore,
\[
    \bar{\z}^T(\s - \bar{\r}) = \bar{\q}^T(\s - \bar{\r}) + t\a^T(\s - \bar{\r}) = \bar{\q}^T(\s - \bar{\r}) \leq 0,
\]
because $\bar{\q} \in \mN_{[0,\gamma]^n}(\bar{\r})$ and $\s \in \Omega \subseteq [0,\gamma]^n$. Hence, $\bar{\z} \in \mN_{\Omega}(\bar{\r})$, and
\begin{equation} \label{eqn:approx-subg}
    H\w + \mu\bar{\r} - \e + \bar{\z} = \bz.
\end{equation}
Let $Q := HH^T + \mu I$. Recall from the proof of Theorem~\ref{thm:dist-to-opt-dual} that for $\z := \nabla g(\y) + L(\y - \r) \in \mN_{\Omega}(\r)$ and $\q := Q\r - \e + \z$,
\[
    \|\q\| \leq \mu\epsilon.
\]
Moreover, $\w = H^T\r$, so $\q = H\w + \mu\r - \e + \z$. Subtracting \eqref{eqn:approx-subg} from this quantity gives
\[
    \q = \mu(\r - \bar{\r}) + \z - \bar{\z}.
\]
By monotonicity of the normal cone,
\begin{align*}
\|\r- \bar{\r}\|\cdot\|\q\| &\geq (\r - \bar{\r})^T\q\\
                            &= \mu\|\r - \bar{\r}\|^2 + (\r - \bar{\r})^T(\z - \bar{\z})\\
                            &\geq \mu\|\r - \bar{\r}\|^2.
\end{align*}
Thus, $\|\r - \bar{\r}\| \leq \|\q\|/\mu \leq \epsilon$. Hence,
\[
    \bar{r}_i \leq r_i + \epsilon < \gamma,
\]
and we know
\[
    \bar{r}_i = \proj_{[0,\gamma]}\left(\frac{1-\theta_i(\w,t)}{\mu}\right).
\]
It follows that $1 - \theta_i(\w,t) < \mu\gamma$, i.e. $\theta_i(\w,t) > 1 - \mu\gamma$.

For the optimal classifier, Theorem~\ref{thm:Csvm-dist-to-opt-dual} gives $r_i^* \leq r_i+\epsilon < \gamma$. Let $(\w^*,t^*)$ be any optimal primal classifier. By Theorem~\ref{thm:Csvm-dual-psi-relation},
\[
    r_i^* = \proj_{[0,\gamma]}\left(\frac{1-\theta_i(\w^*,t^*)}{\mu}\right),
\]
but $r_i^* < \gamma$, so
\[
\frac{1-\theta_i(\w^*,t^*)}{\mu} < \gamma \iff \theta_i(\w^*,t^*) > 1 - \mu\gamma.
\]
\end{proof}

\end{document}
